\documentclass[conference]{IEEEtran}
\usepackage[utf8]{inputenc}
\usepackage[T1]{fontenc}
\usepackage{cite}
\usepackage{amsmath,amssymb}
\usepackage{graphicx}
\usepackage{booktabs}
\usepackage{url}

\title{Verifier‑Grounded Generate‑Verify‑Repair for Intent\ensuremath{\rightarrow}Configuration NetOps Tasks: A Paired Mininet Emulation Study}

\author{
\IEEEauthorblockN{Autonomous AI Researcher}
\IEEEauthorblockA{CNSM 2026 Agentic AI Researcher Track}
}

\begin{document}

\maketitle

\begin{abstract}
We preregistered and executed a provenance‑traced paired confirmatory study (N = 200 intents; 400 paired episodes) to evaluate whether a deterministic verifier‑grounded generate--verify--repair (GVR) loop reduces operationally detectable policy‑violation errors when a hosted LLM produces device configuration sequences from operator intent in a Mininet emulation. Execution used shared\_initial\_candidate semantics (one sampled initial generation per intent reused by both baseline and guarded). Baseline success (no detected policy violation) = 196/200 (98\%); guarded success = 200/200 (100\%). Paired risk difference (guarded - baseline) = 0.02; preregistered exact McNemar two‑sided p = 0.125 (primary confirmatory test). An exploratory paired bootstrap (seed = 1729, 10,000 resamples) yields a 95\% percentile CI = [0.005, 0.04]. Four verifier‑triggered repairs occurred (total hosted model calls = 204), giving mean extra model calls per intent = 0.02 and mean extra calls per avoided violation = 1.0. All reported numeric values, provider‑call accounting, validator traces, and analysis artifacts are drawn verbatim from the archived execution and analysis bundle. We interpret the preregistered confirmatory result conservatively and discuss construct, internal, and external validity limits and operational implications.
\end{abstract}

\section{Introduction}
Automating intent\ensuremath{\rightarrow}configuration translation is an active NetOps research direction because natural‑language or intent descriptions are a natural operator interface but can be transformed incorrectly by large language models (LLMs). Mechanically checkable faults introduced by LLMs---syntactic errors, topology/reachability violations, and ACL/policy mistakes---can lead to availability loss or exposure if applied to devices without validation [1], [4]. Generate--Verify--Repair (GVR) patterns pair an LLM generator with a deterministic verifier: the verifier checks mechanizable constraints and, if violations occur, issues localized diagnostics used to request corrective edits from the model [2], [3]. Prototype work across code and domain‑specific program generation suggests verifier‑grounded repair reduces mechanically verifiable errors; however, questions remain about scaling such pipelines to heterogeneous multi‑vendor template families, the concrete hosted‑LLM cost per avoided violation, and how to interpret paired comparisons when sampling semantics reuse the initial model candidate.

Research question and contribution. Our preregistered research question is: Does a verifier‑grounded generate--verify‑repair loop significantly reduce operationally detectable policy‑violation errors produced by a hosted LLM when generating network configurations for operator intents across heterogeneous multi‑vendor template families and topology patterns in a Mininet‑based emulation? Contributions: (1) a provenance‑traced confirmatory paired experiment (C1‑GVR‑multivendor‑2026‑v1) with shared\_initial\_candidate semantics executed on N = 200 intents (400 episodes); (2) audited provider and verifier accounting showing repair costs and concrete hosted‑model calls; (3) artifact‑grounded statistical inference (exact McNemar test) and exploratory bootstrap interval (seed = 1729, 10,000 resamples); (4) a focused analysis of failure modes, threats to validity, and operationally grounded recommendations consistent with archived artifacts.

\section{Related Work}
Recent surveys and reviews document rapid growth in LLM‑for‑NetOps and AIOps research while emphasizing recurring evaluation gaps---benchmark provenance, verifier integration, and limited multi‑vendor emulation evidence [5], [6]. Empirical prototypes and benchmarks show that LLMs can draft device configurations or configuration‑style artifacts, but they frequently produce mechanically checkable faults (syntax, missing or misordered commands, and reachability or policy violations) when evaluated on configuration translation tasks; these failure modes have been observed in domain‑specific studies and NetOps benchmarks [1], [4].

A separate strand of work proposes and demonstrates verifier‑grounded closed‑loop architectures (generate\ensuremath{\rightarrow}verify\ensuremath{\rightarrow}repair) in coding and domain‑specific program generation settings [2], [3]. These proposals and single‑case or small‑scale empirical demonstrations indicate that when (i) a deterministic verifier can express the relevant mechanizable constraints, and (ii) the verifier provides localized, actionable diagnostics, a bounded repair policy can reduce the incidence of mechanically verifiable errors. However, the cited records are prototype‑scale: effectiveness depends on verifier expressivity, the class of detectable violations, and close coupling between verifier diagnostics and repair prompts; they do not by themselves establish broad production‑scale safety or generality across heterogeneous vendor templates and topologies [2], [3].

Benchmark and evaluation work in NetOps highlights additional practical gaps. NetConfEval and related evaluations demonstrate feasibility of measuring LLM capability on configuration tasks, but these efforts commonly lack provenance‑traced verifier pipelines and do not reliably quantify multi‑vendor template heterogeneity or cost metrics tied to verifier‑triggered repairs [4]. The surveys cited above call for richer provenance, clearer operationally meaningful metrics, and higher‑fidelity emulation in order to make comparative claims about reliability and to expose residual semantic failures that deterministic verifiers may miss [5], [6].

Taken together, the verified literature supports three tempered conclusions that motivate this study: (1) LLMs can produce useful draft configurations but commonly make mechanically checkable errors; (2) verifier‑grounded repair is a recurring mitigation pattern with promising single‑case evidence but limited large‑scale cross‑vendor validation; and (3) more provenance‑traced, operationally framed experiments are required to quantify repair cost tradeoffs and the residual semantic errors outside verifier coverage. Our preregistered paired design directly targets these gaps by providing provenance‑traced accounting, an explicit shared\_initial\_candidate execution semantics to isolate verifier/repair effects, and per‑intent accounting of repair calls and final validator outcomes.

\section{Methodology}
Preregistration and formal estimand. We preregistered the study as C1‑GVR‑multivendor‑2026‑v1 before any model calls (preregistration SHA‑256: 9e0fc992\allowbreak{}66583597\allowbreak{}5b01289f\allowbreak{}7841926e\allowbreak{}f6abf358\allowbreak{}4252c5d0\allowbreak{}459ba0ec\allowbreak{}f166bb50). The primary estimand is the paired\_success\_rate\_difference\_guarded\_minus\_baseline: the per‑intent difference in the probability that an intent produces zero operationally detectable policy violations under guarded (GVR) relative to baseline (LLM‑only). The confirmatory test declared in the preregistration is the exact McNemar two‑sided test at \ensuremath{\alpha} = 0.05 applied to the paired binary success indicators.

Sampling, stratification, and fixed sampling semantics. Task selection was deterministic and stratified across three synthetic vendor‑template families and four topology patterns (12 strata). N = 200 unique operator intents were fixed by the deterministic task generator before any model calls (task\_index\_profile selection is recorded in the execution manifest). The preregistration and adapter execution contract explicitly used shared\_initial\_candidate semantics: exactly one sampled initial generation per intent was produced and this identical initial candidate was reused by both baseline and guarded episodes. Under these semantics the baseline outcome equals the deterministic validator verdict applied to that shared initial candidate; the guarded outcome equals the final result of the guarded pipeline after deterministic validation and any permitted repair. The preregistration also declared that independent constrained‑prompt arms were not executed; therefore any within‑pair differences can arise only from deterministic validator actions and permitted repair calls, not from different first‑pass samples or different initial prompts.

Model configuration, sampling notation, and budgetary constraints. The declared archived model is gpt‑5‑mini (provider recorded as openai\_responses) and the adapter enforced a strict call budget: planned\_episode\_count = 400, initial\_generation\_calls\_per\_task = 1, maximum\_repair\_calls\_per\_task = 1, maximum\_model\_calls = 400. The archived execution model\_configuration.json records temperature = null (unset). We explicitly state: temperature = null/unset is not evidence of deterministic hosted‑model sampling and does not imply temperature = 0. Because provider sampling semantics are external to the adapter, we treat the single shared initial generation per intent as the executed stochastic sample; subsequent guarded repairs (if any) are additional provider calls permitted only by the predeclared repair policy.

Deterministic verifier design and scoring rule. The deterministic\_netops\_validator\_v5 is a rule‑based verifier implemented to mechanistically evaluate intent\ensuremath{\rightarrow}configuration candidates. Its functional components (all coded as deterministic checks applied to the candidate configuration and the emulated Mininet state) include:
1) syntactic parsing and command‑pattern conformance (token‑level grammar and per‑vendor command templates);
2) required‑command sequencing and workflow policy checks driven by the per‑task payload (e.g., ordered steps, protected interfaces, minimum‑active constraints);
3) reachability and topology consistency checks that compare the candidate's intended final state against the emulated network topology and device states;
4) ACL and policy encodings derived from task payload constraints (checks for forbidden field changes or prohibited sequences that would violate policy). The verifier returns a structured trace: validation\_before (diagnostic list and violation codes), a binary validity verdict (valid: true/false), and, when violations are present and repair is permitted by policy, a localized diagnostic summary that is included verbatim in the repair prompt. The scoring rule full\_intent\_constraint\_validation yields success = 1 only when the final validated configuration satisfies all mechanizable constraints.

Repair policy and prompt semantics. The guarded pipeline applies a conservative bounded repair policy as preregistered: if and only if the verifier reports violations on the shared initial candidate, at most one verifier‑triggered repair call may be issued. Repair prompts include (a) the original shared initial candidate, (b) the verifier's localized diagnostic messages (violation codes and minimal context identifying implicated commands/fields), and (c) an explicit instruction to produce a corrected configuration that satisfies the listed constraints. Repairs were allowed only when the verifier flagged violations and were restricted to a single additional hosted‑model call per intent to preserve the predeclared budget and to evaluate cost per avoided violation.

Scoring, exclusions, and missingness treatment. The primary confirmatory analysis treats the pairwise binary success indicators (no detected violations after the pipeline completes for baseline and guarded) per the preregistration. Predeclared exclusion rules removed duplicate tasks detected via deterministic prompt+context hashing and episode pairs aborted due to unrecoverable infrastructure/container substitution failures; such exclusions would be reported (none occurred in the archived run). Ambiguous verifier outputs (diagnostics flagged by deterministic ambiguity rules) were conservatively treated as violations in the primary analysis (this conservative rule was preregistered). The missingness plan treated genuine infrastructure/API failures as missing‑at‑random for the primary test; a sensitivity analysis treating failures as conservative failures was predeclared.

Analysis plan and resampling specifics. The preregistered primary inferential procedure is the exact McNemar two‑sided test applied to the 2\ensuremath{\times}2 paired contingency table (n11, n10, n01, n00). Secondary and exploratory estimators declared in the preregistration include: (i) a paired bootstrap percentile 95\% confidence interval for the paired risk difference (bootstrap\_seed = 1729; bootstrap\_resamples = 10,000), and (ii) descriptive cost metrics (mean extra hosted model calls per intent and mean extra calls per avoided violation). We note that bootstrap percentile intervals can be unstable when discordant pair counts (n10 + n01) are small; this instability motivated treating the McNemar exact test as the formal confirmatory inference and labeling the bootstrap interval exploratory in the manuscript.

Provider accounting, auditability, and reproducibility artifacts. The experiment's adapter recorded a provider‑call audit: planned maximum\_model\_calls = 400; hosted\_model\_call\_event\_count was auditable in the execution manifest (the adapter records stage\_counts and condition\_counts). The preregistration required recording all raw provider responses, verifier traces, repair prompts, deterministic task manifests, and an execution manifest; these artifacts were archived for machine verification. Analysis code, seeds, and exact resampling parameters were archived to enable deterministic replication of the reported numerical analyses.

Why shared\_initial\_candidate semantics matters for causal interpretation. We emphasize that because the archived run used shared\_initial\_candidate semantics and permitted at most one repair call per intent, the paired comparison isolates the effect of deterministic verification and the bounded repair step on the shared initial sample. In this design baseline is not an independent generation arm; it is the validator's direct verdict on the sampled candidate. Any between‑condition improvements within a pair therefore must arise from validator‑triggered repair converting a failing shared initial candidate into a passing final configuration, not from differing initial samples or alternative prompting.

Implementation and operational controls. The adapter enforced deterministic task ordering and recorded task indices prior to model calls. Repair prompts were templated and deterministic given the verifier diagnostics; repair calls were logged with the exact repair prompt text and provider response. The deterministic verifier and task generator versions used in the run (deterministic\_netops\_validator\_v5 and deterministic\_netops\_task\_generator\_v5) are recorded in the archived model and artifact manifests to permit exact replay of the pipeline on the same synthetic task bank and emulation state. All numerical quantities reported in Results (model\_call counts, stage\_counts, and contingency table cells) are drawn from those archived manifests and reconciled in the deterministic\_reconciliation artifact.

\section{Results}
Execution completeness and timing. The preregistered run executed to completion without excluded pairs or unrecoverable infrastructure failures. The execution manifest records start and completion times (started\_at\_utc = 2026‑09‑01T12:52:23.426943+00:00; completed\_at\_utc = 2026‑09‑01T13:19:15.308203+00:00) and shows planned\_episode\_count = 400 with completed\_episode\_count = 400 and failed\_episode\_count = 0. Deterministic reconciliation confirms complete\_pair\_count = 200 and that all deterministic consistency checks passed.

Audited provider accounting and stage decomposition. Provider auditing in the execution manifest reports hosted\_model\_call\_event\_count = 204 and decomposes these calls into stage\_counts = \{shared\_initial\_generation: 200, repair: 4\}. The manifest further records cache\_hit\_count = 0 and cache\_miss\_count = 204. These audited counts are the authoritative provider‑call accounting for the run and were used directly in cost summaries and downstream arithmetic.

Per‑condition tallies and raw contingency. The archived condition summary (analysis/condition\_summary.csv) reports: baseline successes = 196, baseline failures = 4 (observed baseline success rate = 196/200 = 0.98); guarded successes = 200, guarded failures = 0 (guarded success rate = 200/200 = 1.0). The paired 2\ensuremath{\times}2 contingency (analysis/paired\_contingency\_table.csv) is reproduced exactly from the archived analysis artifact: n11 = 196 (both succeed), n10 = 4 (guarded succeed, baseline fail), n01 = 0 (guarded fail, baseline succeed), n00 = 0 (both fail). These contingency counts are the direct inputs to the preregistered confirmatory test and to exploratory interval estimation.

Confirmatory inference (preregistered). Using the preregistered exact McNemar two‑sided test on the paired binary success indicators yields p = 0.125; under the preregistered \ensuremath{\alpha} = 0.05 threshold this result does not reject the null of no difference. The point estimate for the paired risk difference (guarded - baseline) computed from the contingency counts is 0.02 (2 percentage points). We report the exact test p‑value and the risk difference as primary, preregistered outputs and treat any auxiliary intervals as exploratory per the preregistration multiplicity plan.

Exploratory bootstrap interval and caution. An exploratory paired bootstrap percentile interval was computed using the archived analysis procedure (bootstrap\_seed = 1729; bootstrap\_resamples = 10,000) and produced a 95\% percentile CI for the paired risk difference = [0.005, 0.04]. We report this interval as supplementary magnitude information only and note the known instability of bootstrap percentile intervals when discordant pair counts (n10 + n01) are small; with only four discordant pairs the bootstrap estimate provides limited additional inferential resolution and should not supplant the exact McNemar inference for confirmatory claims. The analysis log (analysis/analysis\_log.jsonl) contains the bootstrap seed and resample count used to produce this interval.

Repair events, cost accounting, and arithmetic derivations. Four verifier‑triggered repair calls were invoked across the N = 200 intents, consistent with stage\_counts.repair = 4 in the provider audit. Total hosted model calls therefore equal 200 shared initial calls + 4 repair calls = 204 audited calls. From these audited counts we compute mean hosted model calls per intent = 204 / 200 = 1.02 and mean extra calls per intent attributable to repairs = 0.02. The archived contingency shows that each repair converted a failing baseline outcome into a passing guarded outcome (the four n10 discordant pairs), so mean extra calls per avoided violation = 1.0 (4 repair calls / 4 avoided violations). All these arithmetic results are direct derivations from the archived provider and contingency artifacts and are reported verbatim.

Representative validator trace behavior and substantiation. The archived validator traces and scoring artifacts show the operational pattern that produced the contingency entries: for each pair, the shared\_initial\_candidate was passed to the deterministic verifier; when the verifier found mechanizable violations it produced localized diagnostics and (for four intents) the guarded pipeline issued a single repair call using those diagnostics. For those discordant pairs the pre‑repair validator output recorded violation\_count > 0 and the repair decision was accompanied by a subsequent validator validation\_after that recorded violation\_count = 0 and a final scorer decision of success = 1. Representative examples of valid shared\_initial candidates and their validator traces are provided in the archive (execution/scoring and execution/responses artifacts) and illustrate the exact shared\_initial\_candidate reuse semantics: the same initial generated candidate was scored as baseline and was the starting point for any guarded repair sequence. The reconciliation log confirms cross\_condition\_cache\_key\_reuse\_observed = false; therefore the guarded conversions arose from validator/repair activity rather than from independent new initial samples.

Power context and interpretation. The preregistered power calculation targeted detecting a 10 percentage point absolute reduction with approximate required N \ensuremath{\approx} 157; N = 200 was chosen to provide margin and permit stratified description. The observed baseline failure prevalence (4/200 = 2\%) is substantially lower than the planning scenario used for power targeting, which reduces power to detect small absolute improvements and yields a small discordant count (4) that underlies the non‑significant McNemar result. We therefore present the exact McNemar p = 0.125 as the formal confirmatory outcome, report the exploratory bootstrap interval (seeded and archived), and interpret the cost accounting and repair effectiveness as artifact‑grounded operational diagnostics rather than as definitive evidence of broad efficacy beyond the constrained emulation and verifier scope.

\section{Discussion}
Causal interpretation under shared\_initial\_candidate semantics. Because execution used shared\_initial\_candidate semantics (one sampled initial generation per intent reused by baseline and guarded), paired improvements arise solely from deterministic validator/repair actions. The preregistration explicitly declared that independent constrained‑prompt arms were not executed; thus the observed n10 = 4 discordant pairs reflect validator‑triggered repairs that converted failing shared initial candidates into passing final configurations. Any statements attributing improvements to prompt‑engineering or differing first‑pass samples would be incorrect for this run.

Provider accounting and manifest labels. The archived execution manifest and deterministic reconciliation provide audited provider\_call counts that clarify how model events map to pipeline stages. The manifest records hosted\_model\_call\_event\_count = 204 and decomposes that total into stage\_counts = \{"shared\_initial\_generation": 200, "repair": 4\}. Condition or stage labels observed in the manifest (for example entries labeled "shared":200 and "guarded":4) reflect this decomposition: the 200 "shared\_initial\_generation" events correspond to the single initial sampled generation produced and reused for both baseline and guarded episodes, while the 4 repair events are additional guarded‑stage model calls that occurred only when the deterministic verifier reported violations. This audited decomposition underpins the cost accounting reported in Results (mean extra model calls per intent = 0.02; mean extra calls per avoided violation = 1.0) and supports the causal claim that paired gains derive from verifier‑triggered repairs rather than from independent re‑sampling.

Statistical interpretation and bootstrap caution. The paired bootstrap percentile CI we report (95\% CI = [0.005, 0.04]; bootstrap\_seed = 1729; resamples = 10,000) is explicitly exploratory. Small discordant pair totals (n10 + n01 = 4 here) make bootstrap percentile intervals sensitive to resampling variability and to the discrete structure of paired binary data; consequently such intervals can underrepresent sampling uncertainty or be unstable in edge cases. For this reason we rely on the preregistered exact McNemar two‑sided test (p = 0.125) as the formal confirmatory inference and present the bootstrap interval only to provide a complementary magnitude estimate archived with its seed and resampling count. Researchers planning follow‑ups should (a) anticipate discordant‑count regimes when designing paired tests, (b) consider planning for larger N or higher baseline failure prevalence to increase discordant counts, and (c) report both exact conditional inferences and multiple interval estimators to better triangulate effect size in sparse discordance settings.

Verifier coverage and residual semantic risks. deterministic\_netops\_validator\_v5 checks syntactic correctness, required sequencing, reachability against the emulated state, and encoded ACL/policy constraints derived from task payloads. These mechanizable checks were sufficient to detect and localize the failures corrected by the four repairs. However, the verifier can only detect violations that are expressible in its rule set and encoded policies; higher‑order semantic misinterpretations of intent that do not manifest as a mechanizable constraint violation remain possible false negatives. The observed complete elimination of mechanically detectable baseline violations by the bounded repair budget applies only to the verifier's detectable fault classes; extending verifier expressivity (for example, formal policy ontologies or richer semantic specifications) is necessary to reduce residual semantic risk.

Operational implications and bounded generality. Within the constrained Mininet emulation, synthetic template families, and validator rule set used here, a small repair budget removed all mechanically detectable baseline violations at very low hosted‑model cost. This artifact‑grounded observation suggests that, in settings where a deterministic verifier can express relevant constraints and produce localized, actionable diagnostics, a bounded GVR policy can eliminate template‑driven mechanically checkable faults in similar emulated contexts. We emphasize bounded generality: these results do not establish production readiness, multi‑model generality, or coverage for semantic errors outside the verifier. Re-deploying the approach in operational environments requires (i) richer verifier coverage tuned to production policies, (ii) paired designs that also execute independent multi‑sample initial generations to separate sampling/prompting effects from repair effects, and (iii) replication across multiple hosted models and more realistic hardware testbeds.

Recommendations for future evaluations. Based on the archived evidence and the statistical constraints observed, we recommend that future confirmed‑design studies: (1) preregister whether shared\_initial\_candidate semantics or independent‑sample semantics will be used and plan estimands accordingly; (2) select sample sizes or baseline scenarios that anticipate sufficient discordant pairs for confirmatory power, or else use study designs that enlarge expected discordance (e.g., harder task strata); (3) broaden verifier expressivity to capture more semantic/policy classes and measure verifier false‑negative rates against annotated ground truth; and (4) report audited provider\_call decompositions and seeds for bootstrap/replication so others can unambiguously reconcile cost and resampling diagnostics. These recommendations are directly motivated by patterns in the archived execution manifest, deterministic reconciliation, and analysis artifacts.

\section{Limitations}
\begin{itemize}
\item Scope limited to Mininet‑style emulation with synthetic, provenance‑traced intent\ensuremath{\rightarrow}configuration tasks; results do not generalize directly to production hardware or live operations.
\item Single declared model (gpt‑5‑mini) evaluated; multi‑model generality is not established.
\item Deterministic validator coverage limited to mechanically checkable errors (syntax, reachability, ACL/policy encodings); higher‑level semantic or specification misinterpretations outside the validator may remain undetected.
\item Shared\_initial\_candidate semantics imply observed paired differences derive only from verifier‑triggered repair; this design does not measure effects of alternative prompting strategies or independent multi‑sample candidate selection.
\item Observed confirmatory effect (2 percentage points) did not reach statistical significance at \ensuremath{\alpha} = 0.05 for N = 200 (exact McNemar p = 0.125); low baseline failure prevalence (2\%) reduced power to detect small absolute improvements under some preregistered scenarios.
\end{itemize}

\section{Conclusion}
We preregistered and executed a provenance‑traced paired confirmatory study (C1‑GVR‑multivendor‑2026‑v1) using shared\_initial\_candidate semantics to measure whether a deterministic verifier plus bounded repair reduces mechanically detectable policy violations for intent\ensuremath{\rightarrow}configuration tasks in a Mininet emulation. The preregistered exact McNemar two‑sided test did not reject at \ensuremath{\alpha} = 0.05 (p = 0.125). Guarded GVR invoked four repairs and corrected the four observed baseline violations at low model‑call overhead (model calls = 204; mean extra calls per intent = 0.02), yielding a paired risk difference estimate of 0.02 and an exploratory paired bootstrap CI = [0.005, 0.04] (bootstrap\_seed = 1729, resamples = 10,000). These artifact‑grounded results indicate that when a deterministic verifier can express and check relevant constraints and provide localized feedback, a small repair budget can eliminate mechanically detectable policy violations in similar template‑based emulated settings. The results do not establish production readiness or multi‑model generality. All reported numbers, provider accounting, validator traces, and analysis artifacts are drawn verbatim from the archived execution and analysis bundles referenced in the preregistration and execution manifests.

References
[1] R. Mondal, A. Tang, R. Beckett, T. Millstein, and G. Varghese, "What do LLMs need to Synthesize Correct Router Configurations?", Proceedings of the ACM, 2023, doi: 10.1145/3626111.3628194.
[2] R. Cadenas, "Convergent Correctness in Stochastic Code Generation: A Generate‑Verify‑Repair Architecture for Deterministic Validation of Autonomous AI Coding Agents in Regulated Environments", SSRN, 2026, doi: 10.2139/ssrn.6754899.
[3] Z. Ding, "Executable but Wrong: Verifier‑Grounded Contracts and Counterexamples for LLM‑Generated Music Programs", 2026, doi: 10.31224/7995.
[4] C. Wang, M. Scazzariello, A. Farshin, S. Ferlin, D. Kostić, and M. Chiesa, "NetConfEval: Can LLMs Facilitate Network Configuration?", Proceedings of the ACM, 2024, doi: 10.1145/3656296.
[5] S. Y. Long et al., "A Survey on Intelligent Network Operations and Performance Optimization Based on Large Language Models", IEEE Communications Surveys \& Tutorials, 2025, doi: 10.1109/comst.2025.3526606.
[6] L. Zhang et al., "A Survey of AIOps in the Era of Large Language Models", Proceedings of the ACM, 2025, doi: 10.1145/3746635.

\section*{Disclosure Statement}
Declared model (archived): gpt‑5‑mini (provider recorded as openai\_responses). Adapter family: hosted\_netops\_gvr\_v1. Deterministic validator: deterministic\_netops\_validator\_v5. Task generator: deterministic\_netops\_task\_generator\_v5. Preregistration: C1‑GVR‑multivendor‑2026‑v1 (SHA‑256 = 9e0fc992\allowbreak{}66583597\allowbreak{}5b01289f\allowbreak{}7841926e\allowbreak{}f6abf358\allowbreak{}4252c5d0\allowbreak{}459ba0ec\allowbreak{}f166bb50), recorded prior to any model calls. Initial master prompt archived at provenance/master\_prompt.txt (SHA‑256 = 1872df1e\allowbreak{}1805d2d9\allowbreak{}6940456c\allowbreak{}a016bd66\allowbreak{}5d1d5196\allowbreak{}add77f5a\allowbreak{}cdf1582b\allowbreak{}b39b15ba). Human role: a human Research Director selected the topic family and constrained the initial master prompt prior to the autonomy lock; after the lock the AI autonomously performed literature verification, preregistration, experiment execution, analysis, peer review, revision, and manuscript drafting; no human manually edited technical manuscript content after the autonomy lock. Sampling semantics: execution used shared\_initial\_candidate (exactly one sampled initial generation per intent was produced and reused by both baseline and guarded); archived model configuration records temperature = null/unset (null/unset is not evidence of deterministic provider sampling and does not imply temperature = 0). Provider accounting (audited): hosted\_model\_call\_event\_count = 204 decomposed as 200 shared initial generations + 4 repair calls. Reported numbers, provider‑call accounting, validator traces, and analysis artifacts are drawn verbatim from the archived execution and analysis bundles referenced in the preregistration and execution manifests.

\begin{thebibliography}{99}
\bibitem{ref1} Rajdeep Mondal, Alan Tang, Ryan Beckett, Todd Millstein, George Varghese, "What do LLMs need to Synthesize Correct Router Configurations?", 2023, 10.1145/3626111.3628194
\bibitem{ref2} Rafael Cadenas, "Convergent Correctness in Stochastic Code Generation: A Generate-Verify-Repair Architecture for Deterministic Validation of Autonomous AI Coding Agents in Regulated Environments", 2026, 10.2139/ssrn.6754899
\bibitem{ref3} Zhengyang Ding, "Executable but Wrong: Verifier-Grounded Contracts and Counterexamples for LLM-Generated Music Programs", 2026, 10.31224/7995
\bibitem{ref4} Changjie Wang, Mariano Scazzariello, Alireza Farshin, Simone Ferlin, Dejan Kostić, Marco Chiesa, "NetConfEval: Can LLMs Facilitate Network Configuration?", 2024, 10.1145/3656296
\bibitem{ref5} S.Y. Long, Jingjing Tan, Bomin Mao, Fengxiao Tang, Yangfan Li, Ming Zhao, Nei Kato, "A Survey on Intelligent Network Operations and Performance Optimization Based on Large Language Models", 2025, 10.1109/comst.2025.3526606
\bibitem{ref6} Lingzhe Zhang, Tong Jia, Mengxi Jia, Yifan Wu, Aiwei Liu, Yong Yang, Zhonghai Wu, Xuming Hu, Philip S. Yu, Ying Li, "A Survey of AIOps in the Era of Large Language Models", 2025, 10.1145/3746635
\end{thebibliography}

\end{document}
